We validate the design by ablating the contrastive loss $\mathcal{L}_{\text{NCE}}$ and the vector-quantized loss $\mathcal{L}_{\text{VQ}}$.
The contrastive loss aligns discrete resource usage decisions with the following continuous trajectories.
The vector-quantized loss induces a codebook supporting diverse maneuver patterns under similar discrete actions and state histories.
We denote the ablations as TART w/o $\mathcal{L}_{\text{NCE}}$ and TART w/o $\mathcal{L}_{\text{VQ}}$.
Removing both eliminates the learned representation and conditioning, and the method becomes close to an HPPO~\cite{fan2019hybrid} learner.

In both Maze Navigation and Air-to-Air Combat, ablating either objective reduces \emph{Success Rate}, whereas the full model consistently performs best.
In \textit{Hard} Maze Navigation, both \textbf{TTG} and \emph{Occupancy Coverage} increase in ablated models, reflecting difficulty with clutter and delayed $\textsc{Penetration}$.
In Air-to-Air Combat environment, both ablations increase \textbf{TTE} and \textbf{SPE} with identical budgets, indicating degraded launch timing and weaker follow-up maneuvers.
Variance across seeds is larger for TART w/o $\mathcal{L}_{\text{NCE}}$ than for TART w/o $\mathcal{L}_{\text{VQ}}$, suggesting that multi-modality without contrastive alignment results in unstable learning with mode drift.
These findings highlight the complementary roles of temporal alignment and maneuver diversity for mastering hybrid actions with limited resources under dynamic conditions.